% ================================================================
%  Préambule du manuel de Coupole (FR et EN), d'après le format
%  maison : Latin Modern, titres sans-serif colorés, encadrés
%  tcolorbox, fontawesome5.  \langue{french} ou \langue{english}
%  est défini par le fichier principal avant \input.
% ================================================================
\documentclass[11pt,a4paper]{article}

\usepackage[T1]{fontenc}
\usepackage[utf8]{inputenc}
\usepackage[\langue]{babel}
\usepackage{lmodern}
\usepackage{textcomp}
\usepackage{amsmath,amssymb}
\usepackage{mathtools}
\usepackage{siunitx}
\sisetup{group-minimum-digits=4,separate-uncertainty=true,per-mode=symbol}
\usepackage{cancel}
\usepackage{microtype}
\DeclareUnicodeCharacter{2248}{\ensuremath{\approx}}
\DeclareUnicodeCharacter{2212}{\ensuremath{-}}
\DeclareUnicodeCharacter{00D7}{\ensuremath{\times}}
\DeclareUnicodeCharacter{2192}{\ensuremath{\rightarrow}}
\DeclareUnicodeCharacter{00A0}{~}
\DeclareUnicodeCharacter{202F}{\,}
\DeclareUnicodeCharacter{2033}{\ensuremath{''}}
\DeclareUnicodeCharacter{2032}{\ensuremath{'}}
\DeclareUnicodeCharacter{2264}{\ensuremath{\le}}
\DeclareUnicodeCharacter{2265}{\ensuremath{\ge}}
\DeclareUnicodeCharacter{2194}{\ensuremath{\leftrightarrow}}
\DeclareUnicodeCharacter{00B2}{\textsuperscript{2}}
\DeclareUnicodeCharacter{2026}{\ldots}
\DeclareUnicodeCharacter{00B0}{\textdegree}

\usepackage[a4paper,margin=2cm,headheight=24pt]{geometry}
\usepackage{graphicx}
\graphicspath{{img/}}
\usepackage[table]{xcolor}
\usepackage[most]{tcolorbox}
\usepackage{booktabs}
\usepackage{tabularx}
\usepackage{array}
\usepackage{multicol}
\usepackage{enumitem}
\usepackage[nobottomtitles*]{titlesec}
\usepackage{needspace}
\usepackage{fancyhdr}
\usepackage{fontawesome5}
\usepackage{tikz}
\usetikzlibrary{arrows.meta,calc,decorations.pathreplacing,patterns,angles,quotes}
\usepackage{pgfplots}
\pgfplotsset{compat=1.18}
\usepgfplotslibrary{fillbetween}
\usepackage{pdfpages}
\usepackage[colorlinks=true,linkcolor=sky,urlcolor=sky,bookmarksopen=true,bookmarksopenlevel=1,
  pdftitle={Coupole},
  pdfauthor={ARP273-ROSE}]{hyperref}
\usepackage{bookmark}
% Les notes posées dans des tableaux (annexes) par \footnotemark/\footnotetext cassent les ancres
% de hyperref : les appels de note ne sont donc pas des liens.
\hypersetup{hyperfootnotes=false}
\usepackage{xurl}
\usepackage{fvextra}
\fvset{breaklines=true,fontsize=\small}

% ---------------------------------------------------------------- Langue des titres d'encadrés
\newcommand{\tr}[2]{\ifdefstring{\langue}{english}{#2}{#1}}

% ---------------------------------------------------------------- Couleurs
\definecolor{sky}{HTML}{1C6DB4}
\definecolor{skyl}{HTML}{E7F0F9}
\definecolor{alert}{HTML}{B5382B}
\definecolor{alertl}{HTML}{FBEAE7}
\definecolor{night}{HTML}{1B2A4A}
\definecolor{ink}{HTML}{1F2430}
\definecolor{summit}{HTML}{6B7A8F}
\definecolor{gold}{HTML}{B8860B}
\definecolor{goldl}{HTML}{FBF3DC}
\definecolor{pine}{HTML}{2C7A55}
\definecolor{pinel}{HTML}{E2F1E9}
\definecolor{prune}{HTML}{5B3F8C}
\definecolor{prunel}{HTML}{EEE8F6}
\AtBeginDocument{\color{ink}}

% ---------------------------------------------------------------- Titres
% \section = chapitre du cours (« Chapitre 1 : ... »), \subsection = I, II, III..., \subsubsection = A, B, C...
\titleformat{\section}[display]
  {\sffamily\color{night}}
  {\large\bfseries\color{alert}\MakeUppercase{\sectionname}~\thesection}{2pt}{\LARGE\bfseries}
  [{\vspace{2pt}\color{night!35}\titlerule[1.3pt]}]
\titlespacing*{\section}{0pt}{6pt}{10pt}
\newcommand{\sectionname}{\nomchapitre}
\renewcommand{\thesubsection}{\arabic{section}.\arabic{subsection}}
\renewcommand{\thesubsubsection}{}
\titleformat{\subsection}
  {\large\bfseries\sffamily\color{sky}}{\thesubsection}{0.6em}{}
\titlespacing*{\subsection}{0pt}{12pt}{5pt}
\titleformat{\subsubsection}
  {\normalsize\bfseries\sffamily\color{alert}}{}{0em}{}
\titlespacing*{\subsubsection}{0pt}{9pt}{3pt}
\setcounter{secnumdepth}{3}
\setcounter{tocdepth}{2}

\usepackage{etoolbox}
\pretocmd{\section}{\clearpage}{}{}
\pretocmd{\subsection}{\needspace{8\baselineskip}}{}{}
\pretocmd{\subsubsection}{\needspace{5\baselineskip}}{}{}
\renewcommand{\bottomtitlespace}{0.15\textheight}
\widowpenalty=10000
\clubpenalty=10000
\raggedbottom

% ---------------------------------------------------------------- En-têtes
\pagestyle{fancy}
\fancyhf{}
\renewcommand{\headrulewidth}{0.4pt}
\renewcommand{\headrule}{\hbox to\headwidth{\color{night!40}\leaders\hrule height \headrulewidth\hfill}}
\fancyhead[L]{\sffamily\small\color{night}\textbf{COUPOLE · \titrecourt}}
\fancyhead[R]{\sffamily\small\color{summit}\partiecourante}
\fancyfoot[L]{\sffamily\scriptsize\color{summit}Coupole \versioncoupole{} · GPL-3.0 · ARP273-ROSE}
\fancyfoot[R]{\sffamily\small\color{summit}\thepage}
\newcommand{\partiecourante}{}
\newcommand{\partie}[1]{\clearpage\renewcommand{\partiecourante}{#1}}


% ---------------------------------------------------------------- Parties
\titleformat{\part}[display]{\centering\sffamily\color{night}}
  {\Large\bfseries\color{alert}PARTIE \thepart}{8pt}{\Huge\bfseries}
\pretocmd{\part}{\clearpage}{}{}
\setcounter{tocdepth}{2}

% ---------------------------------------------------------------- Encadrés du site (une icône par type de section)
\tcbset{maison/.style={enhanced,breakable,boxrule=0.9pt,arc=2.5pt,left=2.5mm,right=2.5mm,top=1.5mm,bottom=1.5mm,
  fonttitle=\bfseries\sffamily\small,coltitle=white,
  before upper={\widowpenalties 4 10000 10000 10000 0 \clubpenalties 4 10000 10000 10000 0}}}
% \typesite{env}{icône}{libellé}{fond}{cadre}
\newcommand{\typesite}[5]{\newtcolorbox{#1}[1][]{maison,colback=#4,colframe=#5,title={#2~~#3\ifx\relax##1\relax\else~·~##1\fi}}%
  \BeforeBeginEnvironment{#1}{\par\needspace{6\baselineskip}}}
\typesite{cdefinition}{\faBook}{\tr{Définition}{Definition}}{skyl}{sky}
\typesite{crappel}{\faUndo}{\tr{Rappel}{Reminder}}{goldl}{gold!85!black}
\typesite{cremarque}{\faInfoCircle}{\tr{Remarque}{Note}}{skyl}{sky!70!black}
\typesite{cattention}{\faExclamationTriangle}{\tr{Attention}{Warning}}{alertl}{alert}
\typesite{cexemple}{\faLightbulb}{\tr{Exemple}{Example}}{pinel}{pine}
\typesite{cexercice}{\faPencil*}{\tr{Exercice}{Exercise}}{white}{pine}
\typesite{cdemonstration}{\faCalculator}{\tr{Démonstration}{Proof}}{prunel}{prune}
\typesite{cconseil}{\faThumbsUp}{\tr{Conseil}{Tip}}{pinel}{pine!80!black}
\typesite{ccomplement}{\faPlusCircle}{\tr{Complément}{Supplement}}{prunel}{prune!80!black}
\typesite{censavoirplus}{\faSearchPlus}{\tr{En savoir plus}{Further reading}}{prunel}{prune!70!black}
\typesite{cactivite}{\faFlask}{\tr{Activité}{Activity}}{goldl}{gold}
\typesite{cqcm}{\faCheckSquare}{\tr{QCM}{Quiz}}{white}{sky}
\typesite{cobjectifs}{\faBullseye}{\tr{Objectifs}{Goals}}{skyl}{night}
\typesite{cprerequis}{\faListUl}{\tr{Prérequis}{Prerequisites}}{skyl}{night}
\typesite{ccadre}{\faStar}{\tr{À retenir}{Key point}}{goldl}{gold}
\newtcolorbox{ccode}{maison,colback=night!4,colframe=night!50,fontupper=\small\ttfamily,title={\faTerminal~~\tr{Commande}{Command}}}

\newcommand{\intertitre}[2][]{\par\needspace{4\baselineskip}\medskip\noindent{\sffamily\bfseries\color{night}#2}\par\smallskip}
\newcommand{\question}[1]{\par\smallskip\noindent{\sffamily\bfseries\color{pine}#1}~}
\newcommand{\auteurs}[1]{\par\noindent{\sffamily\small\color{summit}\faUserEdit~#1}\par\medskip}
\newcommand{\appli}{\textcolor{summit}{\faMousePointer~\textit{(application interactive en ligne)}}}
% \figsite{titre}{image}{légende}{crédit}
\newcommand{\figsite}[4]{\par\needspace{8\baselineskip}\begin{center}\begin{minipage}{\linewidth}\centering
  \if\relax\detokenize{#1}\relax\else{\sffamily\small\bfseries\color{sky}#1}\par\smallskip\fi
  \if\relax\detokenize{#2}\relax{\sffamily\small\color{summit}\faFilm~\textit{animation ou vidéo en ligne, non reproduite}}\par\else#2\par\fi
  \if\relax\detokenize{#3}\relax\else\smallskip{\small #3}\par\fi
  \if\relax\detokenize{#4}\relax\else{\scriptsize\color{summit}Crédit : #4}\par\fi
  \end{minipage}\end{center}}
\DeclareUnicodeCharacter{2009}{\,}
\DeclareUnicodeCharacter{2013}{--}
\DeclareUnicodeCharacter{2014}{---}
\DeclareUnicodeCharacter{2019}{'}
\DeclareUnicodeCharacter{2018}{`}
\DeclareUnicodeCharacter{201C}{``}
\DeclareUnicodeCharacter{201D}{''}
% \ressource{titre}{niveau · format · volume}{url}{description} — guide du portail
\newcommand{\ressource}[4]{\par\needspace{5\baselineskip}\noindent{\sffamily\bfseries\color{night}#1}\\
  {\sffamily\footnotesize\color{summit}#2 · \href{#3}{\faExternalLink*~lien}}\\
  {\small #4}\par\smallskip}
